\chapter{Introduction}

Statistical modeling often begins with the analysis of how a response variable $\bm{y}$ varies with respect to one or several predictors $\bm{x}$. In the classical linear regression setting, we are interested in the conditional distribution $p(\bm{y}\mid\bm{x},\bm{\Theta})$, parameterized by regression coefficients and an error variance. Under exchangeability assumptions, the model is expressed as
\[
\bm{y} = \bm{X}\bm{\beta} + \bm{\epsilon}, \quad \bm{\epsilon} \sim \mathsf{N}_n(\bm{0}, \sigma^2 \bm{I}_n),
\]
with $\bm{X}$ the design matrix, $\bm{\beta}$ the regression coefficients, and $\bm{\Theta}=(\bm{\beta},\sigma^2)$ the parameters of interest \cite{GELMAN14}. While this formulation provides a powerful and interpretable framework, it assumes independence across observations, an assumption often violated in real applications.

\medskip

The hierarchical linear regression model (HLRM) extends this framework by allowing data to be structured into groups or clusters, such as students within schools, patients within hospitals, or repeated measures within individuals. In this setting, observations $y_{ij}$ within group $j$ are modeled as
\[
y_{ij} = \bm{x}_{ij}^\top \bm{\beta}_j + \epsilon_{ij}, \quad \epsilon_{ij} \sim \mathsf{N}(0,\sigma_j^2),
\]
with group-specific regression parameters $\bm{\beta}_j$ themselves drawn from a population distribution. This nested specification captures both within- and between-group variability, avoiding the extremes of complete pooling and no pooling \cite{HOFF,SOSA}. Bayesian hierarchical modeling formalizes this perspective by treating regression coefficients and variance components as random, endowed with prior distributions that propagate uncertainty across levels.

\medskip

Estimation in Bayesian HLRMs traditionally relies on Markov Chain Monte Carlo (MCMC) methods. MCMC is asymptotically exact and produces well-calibrated posteriors, but its computational burden is prohibitive in high-dimensional or large-scale problems \cite{GELMAN14}. The search for scalable alternatives has fueled interest in variational inference (VI), a deterministic optimization-based approach that approximates intractable posterior densities by a tractable family of distributions. The roots of VI lie in variational calculus, with modern developments adapting it to Bayesian computation through exponential family theory and conjugacy structures \cite{Diaconis1979-lb, Wainwright2007-ul},\cite{NIPS2000_77369e37,beal2003variational}. 

\medskip

The core idea of VI is to posit a family of candidate densities $q(\bm{\Theta})$ and to select the member that minimizes the Kullback–Leibler divergence to the posterior. This yields the Evidence Lower Bound (ELBO) as an optimization objective. Early expositions demonstrated its utility in Bayesian networks and latent variable models \cite{Blei2006-kd, Ormerod2010-bc}, while comprehensive reviews stress its computational advantages and statistical limitations \cite{Blei2016VariationalIA, zhang2018advancesvariationalinference}. A key requirement is that the model belongs to the natural exponential family, which enables closed-form coordinate updates and efficient implementations \cite{Diaconis1979-lb, NIPS2000_77369e37}. However, this tractability often comes at the cost of approximation quality: VI typically underestimates posterior variances and may distort dependencies among parameters.

\medskip

Modern advances have sought to address these shortcomings. Structured VI relaxes independence assumptions to improve statistical fidelity \cite{wang22g}, while amortized inference with neural networks enables scalable inference in deep generative models, most notably the variational autoencoder (VAE) framework \cite{VAE}. Black-box and stochastic variants, such as stochastic variational inference (SVI), extend applicability to massive datasets by combining stochastic gradients with Robbins–Monro updates \cite{HOFFMAN, RANGAN}. These developments emphasize the trade-off between scalability and accuracy, situating VI and SVI as attractive yet imperfect surrogates to MCMC.

\medskip

This thesis builds upon this body of work by examining the performance of VI and SVI in the context of hierarchical linear regression. We provide a rigorous derivation of the algorithms under conjugate priors, implement their coordinate ascent and stochastic variants, and benchmark them against MCMC in controlled simulations and real-world data. Our contribution is threefold:  
(i) a critical comparison of sampling-based and optimization-based inference across models of increasing complexity (LRM, HLRM, Clustering HLRM);  
(ii) a detailed computational implementation, making variational methods accessible for hierarchical regression analysis; and  
(iii) empirical evidence clarifying when VI and SVI serve as practical substitutes for MCMC, and when their limitations in posterior approximation become decisive.  
In doing so, the thesis not only contributes to the methodological literature by bridging statistical and computational perspectives on variational inference, but also plays a pedagogical role by elucidating its theoretical underpinnings and practical performance.

\medskip

The structure of this thesis is as follows. Chapter 2 introduces the theoretical background, covering the foundations of Bayesian inference, the role of Markov Chain Monte Carlo methods, and the principles of variational inference, including the Kullback–Leibler divergence, the Evidence Lower Bound (ELBO), and both coordinate ascent and stochastic variational approaches. Chapter 3 presents the methodological framework, detailing the formulation of the Linear Regression Model (LRM), the Hierarchical Linear Regression Model (HLRM), and the Clustering HLRM (CHLRM), together with their prior specifications, posterior inference, and illustrative applications. Chapter 4 discusses the main findings, highlighting the methodological and computational trade-offs observed between MCMC, VI, and SVI, and concludes with implications for future research. Finally, the appendices provide complementary material, including notation conventions, full derivations of MCMC and variational updates, and additional implementation details to ensure reproducibility.
